% MT055 / explanation — canonical editable content.
核要像“某个特征空间里的内积”。取样本\(x_1,\ldots,x_n\)，把每对核值填成矩阵\(K_{ij}=k(x_i,x_j)\)。合法实核要求对称，并且对任意系数向量\(a\)都有\(a^TKa\ge0\)，这叫半正定（positive semidefinite，PSD）。


为什么要检查这个量？如果核确实是内积，就有



\[
a^TKa=\sum_{i,j}a_i a_j\phi(x_i)^T\phi(x_j)
=\left\|\sum_i a_i\phi(x_i)\right\|^2\ge0.
\]


向量的平方长度是各分量的平方和，例如\((v_1,v_2)\)的平方长度为\(v_1^2+v_2^2\)，所以右边不可能负。找到一次负值，就能证明候选函数不是所有实输入上的合法核；无需证明它对每组输入都失败。


\textbf{A合法，但排除零向量。} 定义\(\phi(x)=x/\|x\|\)，就是把向量缩到长度1，方向不变。候选式等于\(\phi(x)^T\phi(z)\)，符合内积结构。\(x=0\)时分母为0，题目的原式没有定义。


\textbf{B在积分有限时合法。} \(z\)可理解为一个潜在状态，\(p(x_i\mid z)\)表示固定这个状态后，输入\(x_i\)的概率或密度。先只看两个样本，简写\(p_i(z)=p(x_i\mid z)\)。


对于一个固定的\(z\)，\(i,j\)各取1或2，四项对应\((1,1)\)、\((1,2)\)、\((2,1)\)、\((2,2)\)；中间两项相同，合并后为



\[
\begin{aligned}
\sum_{i,j=1}^{2}a_i a_j p_i(z)p_j(z)
&=a_1^2p_1(z)^2+2a_1a_2p_1(z)p_2(z)+a_2^2p_2(z)^2\\
&=[a_1p_1(z)+a_2p_2(z)]^2.
\end{aligned}
\]


平方项正是这样出现的。若\(z\)只有\(z_1,z_2\)两个可能状态，概率为\(q_1,q_2\ge0\)、\(q_1+q_2=1\)，令\(t_r=a_1p_1(z_r)+a_2p_2(z_r)\)（\(r=1,2\)），再对状态求和便得到\(q_1t_1^2+q_2t_2^2\ge0\)：每个平方非负，权重也非负。


连续的\(z\)用积分把各位置的贡献累加。更多样本仍可按同样方式展开平方，因此对任意系数向量\(a\)，



\[
a^TKa=\int\left(\sum_i a_i p(x_i\mid z)\right)^2p(z)\,dz\ge0.
\]


这里\(p(z)\)是非负密度，\(dz\)表示累加时的小区间宽度。每一份贡献都非负，总和也非负；积分有限的条件保证核值是可用的有限数。


\textbf{C不合法。} 只选一维点\(x_1=0,x_2=2\)，它们的内积分别为0、0、0、4，按候选规则得到



\[
K=\begin{pmatrix}1&1\\1&0\end{pmatrix}.
\]


取\(a=(1,-2)^T\)。先用\(K\)的每一行与\(a\)相乘：



\[
Ka=\begin{pmatrix}1\times1+1\times(-2)\\1\times1+0\times(-2)\end{pmatrix}
=\begin{pmatrix}-1\\1\end{pmatrix}.
\]


再把\(a\)与这个结果求内积，得到\(a^TKa=1\times(-1)+(-2)\times1=-3<0\)，直接反驳半正定。


\textbf{D不合法。} 单点也要满足\(k(x,x)\ge0\)，因为它对应自己的平方长度。取一维\(x=\sqrt{3\pi/2}\)，则\(\sin(x^2)=\sin(3\pi/2)=-1\)，已经失败。


\textbf{E不合法。} \(\tanh\)是双曲正切函数，可按\(\tanh t=(e^t-e^{-t})/(e^t+e^{-t})\)计算。取一维点1、2，核矩阵为\(K=\begin{pmatrix}\tanh1&\tanh2\\\tanh2&\tanh4\end{pmatrix}\)。沿用前面的检查方法，取\(a=(1,-1)^T\)，得到



\[
a^TKa=\tanh1+\tanh4-2\tanh2
\approx0.7616+0.9993-2(0.9640)=-0.1671<0.
\]


所以E也不满足半正定。A、B成立，需保留A的非零输入条件及B的积分有限条件。


\textbf{出处：} 2023B期中 Q4，原卷第3页（5分）


\textbf{答案依据：} 2023B期中 Q4 随卷／配套参考答案；条件与解释经复核。
